\chapter*{Аннотация}
\addcontentsline{toc}{chapter}{Аннотация}

При помощи измерения теплового эффекта Джоуля-Томсона для углекислого газа получили значения поправок a и b в уравнении Ван-Дер-Ваальса:  
$a  = (1,484 \pm 0,242) \  \frac{H\cdot\text{ м} ^4}{\text{моль} ^2}, b = (8,09 \pm 1,86) \cdot 10^{-4} \frac{\text{м} ^3}{\text{моль}}$, и температуры инверсии $T_{\text{инв}} = \frac{2a}{Rb} = (441,3 \pm 124,3) \ K$. Полученные значения разошлись с табличными $a_{\text{табл}} = 0.365$ $\frac{\text{Н·м}^4}{\text{моль}^2}$, $b_{\text{табл}} = 42.79 \times 10^{-4}$ $\frac{\text{м}^3}{\text{моль}}$, ${T_{\text{инв}}}_{\text{табл}} = 2073 K$ в несколько раз. Это говорит о том, что уравнение Ван-дер-Ваальса не позволяет описать поведение углекислого газа.